\chapter{Depression Screening}
\section*{What Is This Test?} Depression screening uses standardized questions to identify depressive symptoms and suicide risk; it is not a standalone diagnosis.
\section*{Alternative Names} PHQ-2; PHQ-9; mood screening; depression questionnaire.
\section*{Principle} A validated questionnaire measures symptom frequency and functional effect, followed by clinical assessment of duration, impairment, mania, substance use, and medical causes.
\section*{Why This Test Is Ordered} It is used in primary care, pregnancy and postpartum care, chronic disease care, and when mood, sleep, appetite, or interest changes are reported.
\section*{Primary vs Secondary Test} Screening is preliminary; diagnosis requires a clinical interview and assessment for bipolar disorder and other explanations.
\section*{How This Test Is Performed} The patient answers questions about mood, interest, sleep, energy, concentration, appetite, self-worth, and self-harm thoughts.
\section*{What Preparation Is Needed} No preparation. Answer honestly and tell the clinician immediately about thoughts of self-harm.
\section*{Understanding Results} A positive screen indicates need for evaluation, not proof of depression. Any suicide-risk response requires prompt safety assessment.
\section*{Follow-up Tests} Follow-up can include psychiatric assessment, substance-use screening, thyroid or other targeted medical tests, and safety planning.
\section*{References} \begin{enumerate}\item MedlinePlus. Depression Screening.\item U.S. Preventive Services Task Force. Depression and suicide-risk screening recommendations.\end{enumerate}
\clearpage\section*{Understanding Results and Follow-up}
The laboratory report should identify the specimen, method, units, reference interval or decision limit, and whether the result is positive, negative, abnormal, or inconclusive. Reference intervals vary with age, sex, pregnancy, specimen type, assay, and laboratory; a result outside a range is not automatically a diagnosis, and a result within a range does not always exclude disease. Discuss unexpected results with the ordering clinician, who can decide whether repeat or confirmatory testing, treatment, or referral is appropriate. Do not change medicines or delay urgent care based on an isolated result.
\section*{Regulatory Status and Authorization Date}This chapter describes a test category, not a specific commercial product. FDA, CDSCO, EU IVDR, and MHRA approval or registration dates therefore cannot be assigned without the assay manufacturer, product name, instrument, and intended use. For laboratory-developed tests, an individual agency approval date may not exist. See the Preface for the interpretation of regulatory status.
\clearpage
